% tikz-tensors-grid.tex -- the grid layout: one tensor repeated over a square
% lattice, drawn turned by 45 degrees so that the lattice reads as a plane
% seen from above -- a PEPS, a two-dimensional partition function. It is a
% block of the canvas, placed as a stack is -- after what is on the row and
% centered on it -- so it takes part in an equation like any stack, and it
% implements the operations of the abstract layout (tikz-tensors-layout.tex)
% that a lattice of this kind has.
%
% A tensor of the grid has a port toward each neighbour on the lattice --
% se and nw along a row (column i to i+1), sw and ne along a column (row j to
% j+1) -- and one out of the plane, down. Split (split), a site is the two
% halves a decomposition of it leaves, and its ports are their corners.
\tn@kind{grid}{}

% \tngrid[bonds=<style>/<label>]{<name>}{<columns>}{<rows>}{<style>/<label>}
% <columns> by <rows> of the tensor <style>/<label>; the one at column i, row j
% is the node <name>-i-j. `bonds=' puts a tensor of its own on every bond,
% half way along -- the weights of a simple update -- turned with the lattice
% so that the bond meets it as it would on a chain, at a vertex of a diamond
% rather than the middle of a side, with its label upright; the one after
% <name>-i-j along a row is <name>-i-j-a, along a column <name>-i-j-b.
% As on a stack, \tnconnect then draws the bonds and \tnopen{down} the
% indices out of the plane; \tnopen{nw}, {ne}, {se}, {sw} the indices of the
% lattice toward that neighbour that no bond takes, and \tnopen{around} all
% four -- a piece of a larger lattice.
%
% [split] draws every site as the two triangles a singular value
% decomposition of it leaves, as the tensor renormalization group splits a
% lattice: on the sites of one colour of the checkerboard (i+j even) the
% indices toward nw and ne go to one half and those toward sw and se to the
% other, the halves one above the other (<name>-i-j-top, -bottom); on the
% other colour the indices toward nw and sw and toward ne and se, the halves
% side by side (-left, -right). Each half is a triangle whose two corners on
% its flat side are its indices on the lattice and whose apex points at its
% other half, and \tnconnect joins the two, apex to apex. The four halves are
% four tensors -- the two factors of each colour's decomposition -- so the
% type is four, <top>, <bottom>, <left>, <right>, each <type>/<label>, or one
% for all four; a notation colours them, and the triangle is the split's.
% scale=, pitch=, size= and stub= draw it at factors of the notation's
% distances and sizes, as on a stack (tikz-tensors-layout.tex).
\pgfkeys{/tn/grid/.is family, /tn/grid/bonds/.initial=,
  /tn/grid/split/.code={\def\tn@gsplit{1}}}
\pgfkeys{/tn/grid/scale/.forward to=/tn@factor/scale,
  /tn/grid/pitch/.forward to=/tn@factor/pitch,
  /tn/grid/size/.forward to=/tn@factor/size,
  /tn/grid/stub/.forward to=/tn@factor/stub}
% the center of column <#2>, row <#3> of grid <#1>, as \tn@gx, \tn@gy
\def\tn@gxy#1#2#3{%
  \tn@len\tn@gx{\tn@get{#1@gx0}+((#2)-(#3)+\tn@get{#1@rows}-1)*\tn@get{#1@gh}}%
  \tn@len\tn@gy{\tn@get{#1@gy0}+((\tn@get{#1@cols}+\tn@get{#1@rows})/2-(#2)-(#3)+1)*\tn@get{#1@gh}}}
\newcommand{\tngrid}[5][]{%
  \def\tn@gsplit{0}%
  \pgfkeys{/tn/grid/.cd, bonds=, #1}%
  \pgfkeysgetvalue{/tn/grid/bonds}\tn@gbond
  \ifnum\tn@gsplit=1 \ifx\tn@gbond\@empty\else
    \PackageError{tikz-tensors}{\string\tngrid: split and bonds= together; a
      split site has no bond tensors}{}\def\tn@gbond{}\fi\fi
  \tn@newblock{#2}{grid}%
  \tn@scaleblock{#2}\tn@tokens{#2}%
  \tn@set{#2@cols}{#3}\tn@set{#2@rows}{#4}%
  % a split site is two tensors where there was one, so its lattice is wider
  \tn@len\tn@gh{\tn@gpitch*0.7071068*\ifnum\tn@gsplit=1 1.35\else 1\fi}%
  \tn@set{#2@gh}{\tn@gh}%
  \tn@block{2\dimexpr\tn@gap\relax}{2*\tn@gap+(#3+#4-2)/2*\tn@gh+\tn@slot/2+\tn@stub}%
  \tn@len\tn@gxz{\the\tn@cursor+\tn@slot/2}\tn@set{#2@gx0}{\tn@gxz}%
  \tn@set{#2@gy0}{\the\tn@line@y}%
  \tn@entry#5/////\tn@end
  \tn@set{#2@split}{\tn@gsplit}%
  \ifnum\tn@gsplit=1 \tn@gfour{#5}\fi
  \foreach \tn@i in {1,...,#3}{\foreach \tn@j in {1,...,#4}{%
    \ifnum\tn@gsplit=1 \tn@ghalves{#2}{\tn@i}{\tn@j}%
    \else
      \tn@gxy{#2}{\tn@i}{\tn@j}%
      \tn@sized{#2}{1}%
      \expandafter\tn@place\expandafter{\tn@opt}{#2-\tn@i-\tn@j}{\tn@gx}{\tn@gy}{\tn@lab}%
    \fi
    \tn@record{#2-\tn@i-\tn@j}{#2}{\tn@i}{\tn@i}{\tn@j}{\tn@j}}}%
  \ifx\tn@gbond\@empty \tn@set{#2@bonds}{0}%
  \else
    \tn@set{#2@bonds}{1}%
    \expandafter\tn@entry\tn@gbond/////\tn@end
    \foreach \tn@i in {1,...,#3}{\foreach \tn@j in {1,...,#4}{%
      \ifnum\tn@i<#3 \tn@gbondnode{#2}{\tn@i}{\tn@j}{\tn@i+1}{\tn@j}{a}\fi
      \ifnum\tn@j<#4 \tn@gbondnode{#2}{\tn@i}{\tn@j}{\tn@i}{\tn@j+1}{b}\fi}}%
  \fi
  \tn@gxy{#2}{#3}{1}\tn@reach{\tn@gx+\tn@slot/2}%
}
% The four types of a split grid's halves, as \tn@ghs@1 .. 4 (top, bottom,
% left, right): four entries, or one for all.
\def\tn@gfour#1{%
  \global\tn@c=0
  \foreach \tn@ge in {#1}{\global\advance\tn@c by 1
    \expandafter\xdef\csname tn@ghs@\the\tn@c\endcsname{\unexpanded\expandafter{\tn@ge}}}%
  \ifnum\tn@c=1
    \foreach \tn@k in {2,3,4}{\expandafter\xdef\csname tn@ghs@\tn@k\endcsname{\tn@ghs@1}}%
  \else\ifnum\tn@c=4 \else
    \PackageError{tikz-tensors}{\string\tngrid: a split grid takes four types,
      top, bottom, left, right, or one; not \the\tn@c}{}\fi\fi}
\def\tn@ghtype#1{\expandafter\expandafter\expandafter\tn@entry
  \csname tn@ghs@#1\endcsname/////\tn@end}
% The two halves of site (<#2>, <#3>) of grid <#1>: one triangle turned four
% ways -- an equilateral one, so that all four halves are one size -- at 0.42
% of the size of the grid's tensors, each placed by its apex (corner 1), the
% apexes a little apart across the site's center. The site keeps its record, for its ports; a
% half is a tensor of its own, for a label or a frame to find.
\def\tn@ghalves#1#2#3{%
  \tn@gxy{#1}{#2}{#3}%
  \tn@len\tn@gg{1.2mm*\tn@get{#1@t@size}}%
  \pgfmathtruncatemacro\tn@gpar{mod(#2+#3,2)}%
  \ifnum\tn@gpar=0
    \tn@ghtype{1}\tn@ghalf{#1}{#1-#2-#3-top}{180}{\tn@gx}{\tn@gy+\tn@gg}%
    \tn@ghtype{2}\tn@ghalf{#1}{#1-#2-#3-bottom}{0}{\tn@gx}{\tn@gy-\tn@gg}%
  \else
    \tn@ghtype{3}\tn@ghalf{#1}{#1-#2-#3-left}{-90}{\tn@gx-\tn@gg}{\tn@gy}%
    \tn@ghtype{4}\tn@ghalf{#1}{#1-#2-#3-right}{90}{\tn@gx+\tn@gg}{\tn@gy}%
  \fi
  \tn@record{#1-#2-#3-\ifnum\tn@gpar=0 top\else left\fi}{#1}{#2}{#2}{#3}{#3}%
  \tn@record{#1-#2-#3-\ifnum\tn@gpar=0 bottom\else right\fi}{#1}{#2}{#2}{#3}{#3}}
\def\tn@ghalf#1#2#3#4#5{% grid, name, turned by, x, y
  \tn@len\tn@hx{#4}\tn@len\tn@hy{#5}%
  \tn@sized{#1}{0.42}%
  \edef\tn@hopt{tn triangle right, shape border rotate=#3, \tn@opt,
                font=\noexpand\scriptsize, anchor=corner 1}%
  \expandafter\tn@ghnode\expandafter{\tn@hopt}{#2}}
\def\tn@ghnode#1#2{\node[#1] (#2) at (\tn@hx,\tn@hy) {\tn@lab};\tn@tracenode{#2}}

% The port of site (<#2>, <#3>) toward <#4> (nw, ne, se, sw), as \tn@px,
% \tn@py: the site's center, which its tensor covers, or on a split site the
% corner of the half that carries that index.
\def\tn@gdirs#1#2#3\tn@end{\def\tn@gns{#1}\def\tn@gwe{#2}}
% which corner of a half is its index toward a direction; corner 1 is its apex
\def\tn@gcn@top@nw{corner 3}\def\tn@gcn@top@ne{corner 2}
\def\tn@gcn@bottom@sw{corner 2}\def\tn@gcn@bottom@se{corner 3}
\def\tn@gcn@left@nw{corner 2}\def\tn@gcn@left@sw{corner 3}
\def\tn@gcn@right@se{corner 2}\def\tn@gcn@right@ne{corner 3}
\def\tn@gport#1#2#3#4{%
  \ifnum\tn@get{#1@split}=1
    \edef\tn@gd{#4}\expandafter\tn@gdirs\tn@gd\tn@end
    \pgfmathtruncatemacro\tn@gpar{mod(#2+#3,2)}%
    \ifnum\tn@gpar=0
      \if n\tn@gns \def\tn@gh{top}\else \def\tn@gh{bottom}\fi
    \else
      \if w\tn@gwe \def\tn@gh{left}\else \def\tn@gh{right}\fi
    \fi
    \tn@anc{#1-#2-#3-\tn@gh}{\csname tn@gcn@\tn@gh @#4\endcsname}%
    \let\tn@px\tn@ax \let\tn@py\tn@ay
  \else
    \tn@gxy{#1}{#2}{#3}\let\tn@px\tn@gx \let\tn@py\tn@gy
  \fi}

% the tensor on the bond from (#2, #3) to (#4, #5), named <#1>-#2-#3-<#6>
\def\tn@gbondnode#1#2#3#4#5#6{%
  \tn@gxy{#1}{#2}{#3}\let\tn@xa\tn@gx \let\tn@ya\tn@gy
  \tn@gxy{#1}{#4}{#5}%
  \tn@len\tn@mx{(\tn@xa+\tn@gx)/2}\tn@len\tn@my{(\tn@ya+\tn@gy)/2}%
  \tn@sized{#1}{1}\edef\tn@opt{\tn@opt, rotate=45}%
  \expandafter\tn@place\expandafter{\tn@opt}{#1-#2-#3-#6}{\tn@mx}{\tn@my}{}%
  \tn@sized{#1}{1}\edef\tn@opt{\tn@opt, draw=none, fill=none}%
  \expandafter\tn@place\expandafter{\tn@opt}{#1-#2-#3-#6-label}{\tn@mx}{\tn@my}{\tn@lab}%
  \tn@record{#1-#2-#3-#6}{#1}{#2}{#2}{#3}{#3}}

% Every bond of the lattice, port to port -- center to center (a tensor covers
% the index it carries), through the tensor on it if there is one, or corner
% to corner on a split grid, which also joins each site's halves.
\def\tn@grid@connect#1{%
  \ifx\tn@apart\@empty\else
    \PackageError{tikz-tensors}{\string\tnconnect: apart= names layers of a
      stack, which a grid has none of, or the ports of bonds to leave out}{}\fi
  \edef\tn@gc{\tn@get{#1@cols}}\edef\tn@gr{\tn@get{#1@rows}}%
  \foreach \tn@i in {1,...,\tn@gc}{\foreach \tn@j in {1,...,\tn@gr}{%
    \ifnum\tn@i<\tn@gc
      \pgfmathtruncatemacro\tn@in{\tn@i+1}%
      \tn@gbondok{#1-\tn@i-\tn@j}{se}{#1-\tn@in-\tn@j}{nw}%
    \fi
    \ifnum\tn@i<\tn@gc\relax \ifnum\tn@gok=1
      \tn@use{#1-\tn@i-\tn@j}{se}{0}\tn@use{#1-\tn@in-\tn@j}{nw}{0}%
      \tn@gport{#1}{\tn@i}{\tn@j}{se}\let\tn@xa\tn@px \let\tn@ya\tn@py
      \tn@gport{#1}{\tn@in}{\tn@j}{nw}\tn@line{\tn@along}{\tn@xa}{\tn@ya}{\tn@px}{\tn@py}%
    \fi\fi
    \ifnum\tn@j<\tn@gr
      \pgfmathtruncatemacro\tn@jn{\tn@j+1}%
      \tn@gbondok{#1-\tn@i-\tn@j}{sw}{#1-\tn@i-\tn@jn}{ne}%
    \fi
    \ifnum\tn@j<\tn@gr\relax \ifnum\tn@gok=1
      \tn@use{#1-\tn@i-\tn@j}{sw}{0}\tn@use{#1-\tn@i-\tn@jn}{ne}{0}%
      \tn@gport{#1}{\tn@i}{\tn@j}{sw}\let\tn@xa\tn@px \let\tn@ya\tn@py
      \tn@gport{#1}{\tn@i}{\tn@jn}{ne}\tn@line{\tn@along}{\tn@xa}{\tn@ya}{\tn@px}{\tn@py}%
    \fi\fi
    \ifnum\tn@get{#1@split}=1
      \pgfmathtruncatemacro\tn@gpar{mod(\tn@i+\tn@j,2)}%
      \ifnum\tn@gpar=0
        \tn@anc{#1-\tn@i-\tn@j-top}{corner 1}\let\tn@xa\tn@ax \let\tn@ya\tn@ay
        \tn@anc{#1-\tn@i-\tn@j-bottom}{corner 1}%
      \else
        \tn@anc{#1-\tn@i-\tn@j-left}{corner 1}\let\tn@xa\tn@ax \let\tn@ya\tn@ay
        \tn@anc{#1-\tn@i-\tn@j-right}{corner 1}%
      \fi
      \tn@line{\tn@along}{\tn@xa}{\tn@ya}{\tn@ax}{\tn@ay}%
    \fi}}}

% An index out of the plane, drawn straight down from tensor <#2> and ending
% at <#3>, a slot and a stub below its center.
\def\tn@gopen#1#2#3{%
  \tn@use{#2}{down}{0}%
  \tn@anc{#2}{center}\let\tn@xa\tn@ax \let\tn@ya\tn@ay
  \tn@len\tn@yb{\tn@ya-\tn@slot/2-\tn@stub}%
  \tn@line{#1}{\tn@xa}{\tn@ya}{\tn@xa}{\tn@yb}%
  \coordinate (#3) at (\tn@xa,\tn@yb);
  \tn@openlabel{#3}}
% An index of the lattice toward <#5> from site (<#3>, <#4>) of grid <#2>,
% that no bond takes: from its port outward along the lattice, as far as an
% index out of the plane shows below a tensor (slot/2 + stub from a center),
% or a stub from the corner of a half; it ends at <#2>-#3-#4-<#5>.
\def\tn@gside#1#2#3#4#5{%
  \tn@use{#2-#3-#4}{#5}{0}%
  \edef\tn@gd{#5}\expandafter\tn@gdirs\tn@gd\tn@end
  \tn@gport{#2}{#3}{#4}{#5}%
  \ifnum\tn@get{#2@split}=1 \tn@len\tn@gl{\tn@stub*0.7071068}%
  \else \tn@len\tn@gl{(\tn@slot/2+\tn@stub)*0.7071068}\fi
  \tn@len\tn@xb{\tn@px\if w\tn@gwe -\else +\fi\tn@gl}%
  \tn@len\tn@yb{\tn@py\if n\tn@gns +\else -\fi\tn@gl}%
  \tn@line{#1}{\tn@px}{\tn@py}{\tn@xb}{\tn@yb}%
  \coordinate (#2-#3-#4-#5) at (\tn@xb,\tn@yb);
  \tn@reach{\tn@xb}%
  \begingroup\let\tn@dir\tn@gd \tn@openlabel{#2-#3-#4-#5}\endgroup}
% may the bond from <#1> toward <#2> to <#3> toward <#4> be drawn? \tn@gok
\def\tn@gbondok#1#2#3#4{\def\tn@gok{1}%
  \tn@hasleg{#1}{#2}\ifnum\tn@hasl=0 \def\tn@gok{0}\fi
  \tn@hasleg{#3}{#4}\ifnum\tn@hasl=0 \def\tn@gok{0}\fi
  \tn@ifcut{#1}{#2}{0}\ifnum\tn@iscut=1 \def\tn@gok{0}\fi
  \tn@ifcut{#3}{#4}{0}\ifnum\tn@iscut=1 \def\tn@gok{0}\fi}
\def\tn@grid@portkey#1#2#3{\def\tn@pkey{0}}
\def\tn@grid@show#1{%
  \foreach \tn@ss in {nw,ne,se,sw,down}{%
    \tn@hasleg{#1}{\tn@ss}%
    \ifnum\tn@hasl=1 \tn@showport{#1}{\tn@ss}{0}\fi}}
% the directions <#1> names: one of nw, ne, se, sw, or around (all four), as
% \tn@gdl; down is \tn@gdl empty
\def\tn@sides{nw,ne,se,sw}\def\tn@saround{around}
\def\tn@gdirlist{%
  \ifx\tn@dir\tn@sdown \let\tn@gdl\@empty
  \else\ifx\tn@dir\tn@saround \let\tn@gdl\tn@sides
  \else \let\tn@gdl\tn@dir
    \edef\tn@gok{\noexpand\in@{,\tn@dir,}{,nw,ne,se,sw,}}\tn@gok
    \ifin@\else
      \PackageError{tikz-tensors}{\string\tnopen: a grid's indices go down, out of
        the plane, or toward nw, ne, se, sw, or around; not `\tn@dir'}{}%
      \let\tn@gdl\@empty\fi
  \fi\fi}
\def\tn@gnosplit#1{\ifnum\tn@get{#1@split}=1
  \PackageError{tikz-tensors}{\string\tnopen: a split grid has no index out of
    the plane}{}\fi}
\def\tn@grid@opennode#1#2{%
  \tn@gdirlist
  \edef\tn@gb{\tn@get{n@#2@s}}\edef\tn@gi{\tn@get{n@#2@a}}\edef\tn@gj{\tn@get{n@#2@k}}%
  \ifx\tn@gdl\@empty
    \expandafter\tn@gnosplit\expandafter{\tn@gb}%
    \tn@free{#2}{down}{0}{\string\tnopen}\tn@gopen{#1}{#2}{#2-down}%
  \else
    \foreach \tn@gs in \tn@gdl{%
      \ifx\tn@dir\tn@saround
        \tn@ifused{\tn@gb-\tn@gi-\tn@gj}{\tn@gs}{0}%
      \else \tn@free{\tn@gb-\tn@gi-\tn@gj}{\tn@gs}{0}{\string\tnopen}\def\tn@isused{0}\fi
      \ifnum\tn@isused=0 \tn@gside{#1}{\tn@gb}{\tn@gi}{\tn@gj}{\tn@gs}\fi}%
  \fi}
\def\tn@grid@openblock#1#2{%
  \tn@gdirlist
  \edef\tn@gc{\tn@get{#2@cols}}\edef\tn@gr{\tn@get{#2@rows}}%
  \ifx\tn@gdl\@empty \tn@gnosplit{#2}\fi
  \foreach \tn@i in {1,...,\tn@gc}{\foreach \tn@j in {1,...,\tn@gr}{%
    \ifx\tn@gdl\@empty
      \tn@ifused{#2-\tn@i-\tn@j}{down}{0}%
      \ifnum\tn@isused=0 \tn@gopen{#1}{#2-\tn@i-\tn@j}{#2-\tn@i-\tn@j-down}\fi
    \else
      \foreach \tn@gs in \tn@gdl{%
        \tn@ifused{#2-\tn@i-\tn@j}{\tn@gs}{0}\tn@hasleg{#2-\tn@i-\tn@j}{\tn@gs}%
        \ifnum\tn@isused\tn@hasl=01 \tn@gside{#1}{#2}{\tn@i}{\tn@j}{\tn@gs}\fi}%
    \fi}}}
